\chapter{Información General del Proyecto}

\section{Ubicación Geográfica y Fecha de Entrada en Operación}

El proyecto está ubicado en Bucaramanga, departamento de Santander con entrada en operación para 2025.

\begin{table}[H]
    \centering
    \caption{Coordenadas ubicación geográfica del proyecto}
    \begin{tabular}{cc}
    \toprule
    \textbf{Latitud} & \textbf{Longitud} \\
    \midrule
    7,12621563 & -73,11227598 \\
    \bottomrule
    \end{tabular}
\end{table}

\section{Parámetros Técnicos de Inversores y Paneles}

\subsection{Inversores}

\begin{table}[H]
    \centering
    \caption{Parámetros técnicos del inversor SOLIS 60K LV 5G}
    \begin{tabular}{lcc}
    \toprule
    \textbf{Parámetro} & \textbf{Valor} & \textbf{Unidad} \\
    \midrule
    Número de inversores & 2 & UND \\
    Tensión máxima de arranque & 195 & VDC \\
    Tensión MPPT en DC entrada & 180 - 1000 & VDC \\
    Tensión máxima DC entrada & 1100 & VDC \\
    Máxima corriente de entrada & $26 \times 8$ & A DC \\
    Potencia máxima & 69000 & W \\
    Voltaje nominal de salida & 220 & VAC \\
    Frecuencia de salida & 60 & Hz \\
    Corriente de salida & 157,5 & A \\
    Número de fases & 3 & - \\
    Protección anti-isla & Sí & - \\
    Factor de potencia & 0,99 & - \\
    \bottomrule
    \end{tabular}
\end{table}

\subsection{Paneles Fotovoltaicos}

\begin{table}[H]
    \centering
    \caption{Parámetros técnicos de los paneles fotovoltaicos}
    \begin{tabular}{lcc}
    \toprule
    \textbf{Parámetro} & \textbf{Valor} & \textbf{Unidad} \\
    \midrule
    Tipo & \multicolumn{2}{l}{JINKO Monocristalino Mono-facial} \\
    Potencia nominal por panel & 590 & W \\
    Voltaje máximo de salida & 41,05 & VDC \\
    Corriente máxima de salida & 10,83 & A \\
    Eficiencia & 22,84 & \% \\
    Cantidad de paneles & 240 & UND \\
    \bottomrule
    \end{tabular}
\end{table}

\section{Cálculo de Energía Producida}

\subsection{Radiación Solar}

Teniendo en cuenta la radiación diaria promedio mensual y el número de días que conforman los doce meses del año, se puede determinar el valor de la radiación global anual sobre una superficie horizontal para la ubicación de Bucaramanga, Santander. Este valor es igual a 4,7 kWh/m² en promedio.

\subsection{Ausencia de Sombras}

En el caso de la planta en cuestión, no existe sombreado que afecte la generación del sistema fotovoltaico.

\subsection{Cálculo de Generación de Energía}

La producción de energía del sistema se calculó sobre la base de datos climáticos de la PVGIS NSRDB/ERA-Interim del lugar de la instalación.

\begin{table}[H]
    \centering
    \caption{Datos de entrada para estimar la energía generada}
    \begin{tabular}{lc}
    \toprule
    \textbf{Parámetro} & \textbf{Valor} \\
    \midrule
    Ubicación del proyecto & Bucaramanga \\
    Coordenadas (LAT, LONG) & 7.12621563, -73.112275 \\
    Capacidad instalada en DC & 141,6 kWp \\
    Capacidad instalada en AC & 120 kW \\
    Tipo de módulo & Premium \\
    Inclinación & 10° \\
    Azimut & 180 \\
    Pérdidas del sistema & 26\% \\
    Eficiencia del inversor & 98,50\% \\
    Relación DC a AC & 1,18 \\
    \bottomrule
    \end{tabular}
\end{table}

\subsection{Fórmula de Energía Producida}

La energía producida por el sistema sobre una base anual se calcula mediante:

\begin{equation}
E_{p,y} = P_{nom}(d.c) \times Irr \times (1 - \text{Losses}) = 179.756,95 \text{ kWh/día}
\end{equation}

Donde:
\begin{itemize}
    \item $P_{nom}(d.c)$ = Potencia Nominal del sistema: 141,6 kWp
    \item $Irr$ = Radiación anual sobre la superficie de los módulos: 1.715,5 kWh/m²
    \item Losses = Pérdidas de potencia: 26\%
\end{itemize}

Las pérdidas de potencia se dan debido a varios factores como pérdidas por efecto Joule, disminución de la eficiencia de los equipos debido a la temperatura, entre otros.

\section{Parámetros Técnicos del Nodo 3272966}

Los parámetros eléctricos de las líneas del Sistema de Distribución Local (SDL) fueron suministrados por el operador de red, correspondientes a la base utilizada para modelar el sistema de distribución de ESSA.

El nodo 3272966 se conecta al sistema de distribución a través del tramo de línea 3272869 – 3272966, cuyas características eléctricas son:

\begin{table}[H]
    \centering
    \caption{Parámetros de la línea de conexión}
    \begin{tabular}{lc}
    \toprule
    \textbf{Parámetro} & \textbf{Valor} \\
    \midrule
    Conductor tipo & 3F\_15\_CU\_2\_XLPE \\
    Longitud & 0,003 km \\
    Corriente nominal & 0,155 kA \\
    Resistencia de secuencia positiva (R1) & 0,0020 Ω \\
    Reactancia de secuencia positiva (X1) & 0,0013 Ω \\
    Resistencia de secuencia cero (R0) & 0,0033 Ω \\
    Reactancia de secuencia cero (X0) & 0,0012 Ω \\
    \bottomrule
    \end{tabular}
\end{table}

\subsection{Equivalente de Red}

Para efectos de los análisis de cortocircuito, el equivalente de red en la subestación CONUCOS, a un nivel de tensión de 13,8 kV, se caracteriza mediante los siguientes niveles de falla:

\textbf{Falla Trifásica:}
\begin{itemize}
    \item Potencia de cortocircuito ($S_k$): 654,49 MVA
    \item Corriente simétrica de cortocircuito ($I_k$): 27,38 kA
    \item Corriente de cresta ($I_p$): 70,40 kA
\end{itemize}

\textbf{Falla Monofásica a Tierra:}
\begin{itemize}
    \item Potencia de cortocircuito: 239,68 MVA
    \item Corriente simétrica: 30,08 kA
    \item Corriente de cresta: 77,35 kA
    \item Relación R/X del equivalente: 0,068
\end{itemize}

\section{Transformador de Conexión}

\begin{table}[H]
    \centering
    \caption{Parámetros técnicos del transformador Serie 21123832-S}
    \begin{tabular}{lcc}
    \toprule
    \textbf{Parámetro} & \textbf{Valor} & \textbf{Unidad} \\
    \midrule
    Potencia & 150 & kVA \\
    Tensión primario / secundario & 13,2 / 220 & kV / V \\
    Impedancia de cortocircuito (Zcc) & 4,9 & \% \\
    Grupo de conexión & Dyn5 & - \\
    \bottomrule
    \end{tabular}
\end{table}
